\documentclass[11pt]{article}
\usepackage[margin=1in]{geometry}
\usepackage{amsmath,amssymb,amsthm}
\usepackage[hidelinks]{hyperref}
\usepackage{graphicx}

\newtheorem{theorem}{Theorem}
\newtheorem{lemma}{Lemma}
\newtheorem{corollary}{Corollary}
\newtheorem{remark}{Remark}

\newcommand{\HG}{\mathcal H}
\newcommand{\norm}[1]{\lVert #1\rVert}
\newcommand{\ran}{\operatorname{ran}}
\newcommand{\Reals}{\mathbb R}
\newcommand{\Complex}{\mathbb C}

\title{The infinite-dimensional Hilbert version of the relative spear theorem}
\author{FA Banach run \texttt{fa\_banach\_001}}
\date{June 30, 2026}

\begin{document}
\maketitle

\section*{Source question}

The source is Lakshmi Kanta Dey and Subhadip Pal, \emph{On the geometry of
$G$-norm}, arXiv:2603.20796.  Question 1 on PDF page 13 asks:

\begin{quote}
What is the infinite-dimensional version of Theorem 3.12?
\end{quote}

Here Theorem 3.12 is the finite-dimensional Hilbert-space theorem asserting
that a relative spear operator is a partial isometry.  The source question is
shown below.

\begin{center}
\includegraphics[width=0.92\linewidth]{figures/open_problem_crop.png}
\end{center}

\section*{Definitions}

Let $H_1,H_2$ be Hilbert spaces over $\mathbb K=\Reals$ or $\Complex$, and use
the inner product convention linear in the first variable.  Let
$G\in L(H_1,H_2)$ with $\norm{G}=1$.  For $T\in L(H_1,H_2)$, set
\[
  \norm{T}_G
  = \inf_{\delta>0}\sup\{\norm{Tx}:x\in S_{H_1},\ \norm{Gx}>1-\delta\}.
\]
The $G$-numerical radius is
\[
  \nu_G(T)=
  \sup\{|\lambda|:\lambda\in V_G(T)\},
\]
where $V_G(T)$ consists of the limits of $ \langle T x_i,y_i\rangle $ along
nets with $x_i\in S_{H_1}$, $y_i\in S_{H_2}$ and
$\Re\langle Gx_i,y_i\rangle\to 1$.  We call $G$ \emph{relative spear} if
\[
  \nu_G(T)=\norm{T}_G\qquad(T\in L(H_1,H_2)).
\]

\section*{Theorem}

\begin{theorem}\label{thm:classification}
Let $H_1,H_2$ be Hilbert spaces and let $G\in L(H_1,H_2)$ have norm one.
Then $G$ is relative spear if and only if $H_2$ is one-dimensional and $G$ is a
rank-one partial isometry onto $H_2$.

Consequently, if $\norm{\cdot}_G$ is required to be a genuine norm on
$L(H_1,H_2)$, then $G$ is relative spear if and only if $H_1$ and $H_2$ are both
one-dimensional and $G$ is a surjective isometry.
\end{theorem}

\section*{Proof intuition}

The finite-dimensional theorem detects singular values in $(0,1)$ by mapping a
top singular vector into a lower singular direction.  In infinite dimensions the
same obstruction has two forms.  If the top spectral value of $G^*G$ is
isolated, the finite-dimensional argument works with the spectral projection at
$1$.  If spectral values below $1$ accumulate at $1$, one uses two alternating
families of near-top spectral slices and builds a nilpotent shift between them.
That operator still has $G$-norm one because its domain slices approach the top
of the spectrum, but its $G$-numerical radius is bounded by $1/2$.

Once $G^*G$ is forced to be a projection, $G$ is a partial isometry.  A
relative spear partial isometry cannot have missing final directions, because a
test operator can point into the orthogonal complement of the final space and
be invisible to $\nu_G$.  It also cannot have a final space of dimension at
least two, because a two-dimensional nilpotent block has numerical radius
$1/2$.  Thus the only Hilbert relative spear case is one-dimensional in the
codomain.

\section*{Proof}

We first record a Hilbert-space simplification of $V_G$.

\begin{lemma}\label{lem:hilbert-vg}
If $\Re\langle Gx_i,y_i\rangle\to 1$ with
$x_i\in S_{H_1}$ and $y_i\in S_{H_2}$, then
$\norm{Gx_i}\to 1$ and $\norm{y_i-Gx_i}\to0$.  Hence every element of
$V_G(T)$ is a limit of $\langle T x_i,Gx_i\rangle$ along a net with
$\norm{Gx_i}\to1$.
\end{lemma}

\begin{proof}
Since $\norm{Gx_i}\le1$ and $\norm{y_i}=1$, the condition
$\Re\langle Gx_i,y_i\rangle\to1$ forces $\norm{Gx_i}\to1$.  Moreover
\[
  \norm{y_i-Gx_i}^2
  =1+\norm{Gx_i}^2-2\Re\langle Gx_i,y_i\rangle \to0.
\]
Thus
$|\langle T x_i,y_i\rangle-\langle T x_i,Gx_i\rangle|
\le \norm{T}\,\norm{y_i-Gx_i}\to0$.
\end{proof}

Let $A=G^*G$.  Then $A$ is a positive contraction on $H_1$ with $\norm{A}=1$.
We show first that $A$ must be a projection.

Assume, toward a contradiction, that the spectral projection
$E_A((0,1))$ is nonzero.

\medskip
\noindent
\textbf{Case 1: the top is isolated from the rest of the spectrum.}
Suppose that $1$ is isolated from $\sigma(A)\setminus\{1\}$.  Then
$M=E_A(\{1\})H_1$ is nonzero, and there is $\gamma<1$ such that
$\norm{A^{1/2}x}\le \gamma\norm{x}$ on $M^\perp$.  Choose a unit
$u\in M$ and a unit vector $v\in E_A((0,1))H_1$.  Put
$e=Gv/\norm{Gv}$ and define
\[
  T x=\langle x,u\rangle e.
\]
Since $\norm{Gu}=1$ and $\norm{Tu}=1$, we have $\norm{T}_G=1$.  If
$\norm{Gx_i}\to1$, then writing $x_i=P_Mx_i+(I-P_M)x_i$ gives
$\norm{(I-P_M)x_i}\to0$.  Also $Gv\perp G M$, so
\[
  \langle T x_i,Gx_i\rangle
  =\langle x_i,u\rangle \langle e,G(I-P_M)x_i\rangle \to0.
\]
By Lemma \ref{lem:hilbert-vg}, $\nu_G(T)=0$, contradicting relative spear.

\medskip
\noindent
\textbf{Case 2: spectral values below $1$ accumulate at $1$.}
Choose pairwise disjoint Borel sets $I_n,J_n\subset(0,1)$ with nonzero spectral
projections and with $\inf I_n,\inf J_n\to1$.  Pick unit vectors
$u_n\in E_A(I_n)H_1$ and $v_n\in E_A(J_n)H_1$.  The families $(u_n)$ and
$(v_n)$ are mutually orthonormal, and the vectors
$e_n=Gv_n/\norm{Gv_n}$ are orthonormal in $H_2$.  Define
\[
  T x=\sum_{n=1}^\infty \langle x,u_n\rangle e_n .
\]
This is a contraction.  Since $\norm{G u_n}\to1$ and $\norm{T u_n}=1$, it
follows that $\norm{T}_G=1$.

For every $x\in S_{H_1}$,
\[
  \langle T x,Gx\rangle
  =\sum_n \langle x,u_n\rangle\, \langle e_n,Gx\rangle .
\]
The coefficient vector
$(\langle e_n,Gx\rangle)_n$ depends only on the projection of $x$ onto
$\overline{\operatorname{span}}\{v_n:n\ge1\}$: indeed
$\langle e_n,Gx\rangle=\langle A v_n,x\rangle/\norm{Gv_n}$, and the mutually
orthogonal vectors $A v_n/\norm{Gv_n}$ have norm at most one because
$0\le A\le I$.  Hence this coefficient vector has norm at most
$\|P_Vx\|$, and
\[
  |\langle T x,Gx\rangle|
  \le
  \|P_Ux\|\,\|P_Vx\|
  \le \frac12,
\]
where $P_U$ and $P_V$ are the orthogonal projections onto the closed spans of
$(u_n)$ and $(v_n)$.  Lemma \ref{lem:hilbert-vg} gives
$\nu_G(T)\le1/2<1=\norm{T}_G$, again a contradiction.

Both cases are impossible; hence $E_A((0,1))=0$.  Therefore
$A=G^*G$ is an orthogonal projection, and $G$ is a partial isometry.  Let
$M=\ran A$ be the initial space and $N=GM$ be the final space.

If $N\ne H_2$, choose a unit $u\in M$ and a unit $z\in N^\perp$.  The operator
$T x=\langle x,u\rangle z$ satisfies $\norm{T}_G=1$ because $\norm{Gu}=1$.
But $Gx\in N$ for all $x$, so
$\langle Tx,Gx\rangle=0$ for all $x$.  Thus $\nu_G(T)=0$, contradiction.
Consequently $N=H_2$.

If $\dim H_2\ge2$, then $\dim M\ge2$.  Choose orthonormal $u_1,u_2\in M$ and
let $S$ be the rank-one nilpotent on $M$ given by
$S u_1=u_2$ and $S=0$ on $u_1^\perp\cap M$.  Define
$T=G S P_M$, where $P_M$ is the projection onto $M$.  Then $\norm{T}_G=1$.
For all $x\in H_1$,
\[
  \langle Tx,Gx\rangle
  =\langle S P_Mx,P_Mx\rangle .
\]
The numerical radius of this nilpotent block is $1/2$, so
$|\langle Tx,Gx\rangle|\le1/2$ for $x\in S_{H_1}$.  Lemma
\ref{lem:hilbert-vg} yields $\nu_G(T)\le1/2$, contradiction.  Therefore
$\dim H_2=1$.

Since $N=H_2$ and $G$ is a partial isometry, $G$ is a rank-one partial isometry
onto the one-dimensional codomain.

Conversely, suppose $H_2=\mathbb K v$ and
\[
  Gx=\langle x,u\rangle v
\]
for unit vectors $u\in H_1$ and $v\in H_2$.  Let $T\in L(H_1,H_2)$, so
$Tx=f(x)v$ for some bounded functional $f$ on $H_1$.  If
$x=\alpha u+h$ with $h\perp u$ and $|\alpha|>1-\delta$, then
$\norm{h}\le (2\delta)^{1/2}$.  It follows that
\[
  \norm{T}_G=\lim_{\delta\downarrow0}
  \sup_{\substack{x\in S_{H_1}\\ |\langle x,u\rangle|>1-\delta}} |f(x)|
  =|f(u)|.
\]
On the other hand, if
$\Re\langle Gx_i,y_i\rangle\to1$, then $y_i=\beta_i v$ with
$|\beta_i|=1$, $\langle x_i,u\rangle\overline{\beta_i}\to1$, and
$x_i-\langle x_i,u\rangle u\to0$.  Hence
\[
  \langle Tx_i,y_i\rangle
  =f(x_i)\overline{\beta_i}\to f(u).
\]
Thus $\nu_G(T)=|f(u)|=\norm{T}_G$ for every $T$, so $G$ is relative spear.

Finally, if $\norm{\cdot}_G$ is a genuine norm on $L(H_1,H_2)$ and
$Gx=\langle x,u\rangle v$, then $u^\perp=\{0\}$; otherwise a nonzero operator
supported on $u^\perp$ has $G$-seminorm zero.  Hence $H_1$ is also
one-dimensional.  The converse in the one-dimensional case is immediate.
\qed

\section*{Answer to the Source Question}

The infinite-dimensional Hilbert version of the source theorem is not merely
``relative spear implies partial isometry.''  The exact classification is:
relative spear Hilbert operators are precisely the rank-one partial isometries
onto a one-dimensional codomain.  With the source paper's standing convention
that the $G$-norm is a norm on all of $L(H_1,H_2)$, the only Hilbert examples
are one-dimensional surjective isometries.

\section*{Novelty Check and Limitations}

Run-local indexes were searched for arXiv:2603.20796 and for the terms
\texttt{G-norm}, \texttt{relative spear}, \texttt{partial isometry}, and
Hilbert-space variants.  No prior run packet was found.  A targeted web search
for the exact title and relative-spear/Hilbert/partial-isometry phrases did not
locate an obvious later answer.

This packet answers Question 1.  It does not settle the source paper's broader
Banach-space Questions 2 and 3.

\section*{References}

\begin{enumerate}
\item Lakshmi Kanta Dey and Subhadip Pal, \emph{On the geometry of
$G$-norm}, arXiv:2603.20796.
\end{enumerate}

\end{document}
